\chapter{부록 --- 도구 설정}
\label{ch:toolchain}

\section{Python 환경}
\label{sec:python-env}

이 수업의 코드는 \texttt{uv} 로 관리한다. \texttt{pip} 나 \texttt{conda} 는 쓰지
않는다.

\begin{lstlisting}[language=bash]
uv venv --python "$(pyenv which python)"   # 가상환경 생성
uv add pandas matplotlib scipy openpyxl    # 의존성 추가
uv run python hw1/code/p1.py               # 실행
\end{lstlisting}

\begin{pitfall}[matplotlib 백엔드]
이 가상환경에는 \texttt{\_tkinter} 가 없다. 파일만 저장할 때는 \texttt{pyplot} 을
import 하기 \emph{전에} 백엔드를 고정해야 한다.

\begin{lstlisting}[language=Python]
import matplotlib
matplotlib.use("Agg")       # 반드시 pyplot import 앞에
import matplotlib.pyplot as plt
\end{lstlisting}
\end{pitfall}

\section{이 책의 빌드}
\label{sec:book-build}

\begin{lstlisting}[language=bash]
cd book_intro_2_ds
make          # latexmk -xelatex -> main.pdf
make clean    # 중간 파일 정리 (PDF 는 남김)
\end{lstlisting}

엔진은 \textbf{XeLaTeX} 이다. 한글 조판에 \texttt{xeCJK} 가 필요하고, 이 기기의
TeX Live 는 basic 스킴이라 \texttt{kotex} 이 없다. \texttt{pdflatex} 으로는
빌드되지 않는다.

\begin{aside}[없는 패키지와 그 대체]
basic 스킴에 \texttt{siunitx}, \texttt{tcolorbox}, \texttt{mdframed},
\texttt{framed}, \texttt{biblatex} 가 없다. 콜아웃 상자는 \texttt{tcolorbox} 대신
\texttt{\textbackslash fcolorbox} 로 직접 만들었다(\texttt{preamble.tex}).
필요하면 \texttt{sudo tlmgr install <패키지>} 로 설치한다.
\end{aside}

\section{그림 재생성}
\label{sec:figure-regen}

책의 그림은 복사본이 아니라 과제 폴더를 직접 참조한다(\texttt{main.tex} 의
\texttt{\textbackslash graphicspath}). 따라서 코드를 다시 돌리면 책도 함께
갱신된다.

\begin{lstlisting}[language=bash]
.venv/bin/python hw1/code/p1_figures.py         # fig1-3
.venv/bin/python hw1/code/p1_stats_figures.py   # q1_fig1-3
.venv/bin/python hw1/00_lecture_charts/main.py  # 강의 차트 21개
\end{lstlisting}
